\section{Finite front-door identification: proofs}
\label{sec:annex-p16-front-door-identification}

Use the finite model, graph and intervention notation of
\cref{def:p16-front-door-data}. The source's mediated composition is
Eq.~(26) of \cite[sec.~3.3.2, p.~116]{Pearl2009CausalOverview}.
The proofs here establish the displayed graph criterion and support
conditions without using a general do-calculus rule.

\begin{annexitem}{annex:p16.front-door-graph-steps}{thm:p16-front-door-graph-steps}{Three graph consequences of the criterion}
\begin{proof}
First, suppose \(t\) were a strict descendant of \(h\). Reverse
a directed path from \(h\) to \(t\). It is a back-door path
from \(t\) to \(h\) with no collider, so it is active given
the empty set. This contradicts the second front-door condition.
Thus \(t\notin\operatorname{de}_G(\{h\})\). A directed path
in a DAG has distinct vertices, so the path convention introduces
no additional case here.

Second, suppose there were a directed path from \(t\) to \(r\)
in \(G_{\overline h}\). It is also a directed path in \(G\),
so the first condition requires it to contain \(h\). Since its
distinct endpoints are \(t,r\), this occurrence of \(h\) is
internal and has an incoming path edge. Such an edge has been deleted
from \(G_{\overline h}\), a contradiction. Therefore no such
directed path survives.

For the third assertion, suppose a back-door path from \(h\) to
\(r\) in \(G_{\overline t}\) is active given the empty set.
It has no collider, because the empty set activates none. Its edges
and collider status are unchanged in \(G\), where it is still a
back-door path. The third front-door condition says that \(\{t\}\)
blocks it in \(G\). As it has no collider, blocking is possible
only because an internal noncollider on it is \(t\).

There is no incoming edge at \(t\) in \(G_{\overline t}\).
Consequently both path edges at this internal vertex point away from
\(t\). Consider the portion from \(t\) toward the endpoint
\(r\). Its first edge points in the direction of travel. At each
subsequent internal vertex, an arriving arrow points into that vertex;
absence of a collider forces the next arrow to point away from it.
Induction along the portion therefore makes it a directed path from
\(t\) to \(r\). It avoids \(h\), which is the other endpoint
of the original simple path. This contradicts the first front-door
condition. There is no active back-door path of the proposed kind,
proving the third assertion.
\end{proof}
\end{annexitem}

\begin{annexitem}{annex:p16.front-door-composition}{thm:p16-front-door-composition}{Composition proved within the intervened models}
\begin{proof}
Fix \(x\in S_t\) and \(m\in S_h\). First intervene at the
mediator, obtaining the model \(\mathcal M_{h=m}\) with graph
\(G_{\overline h}\). By
\cref{thm:p16-front-door-graph-steps}, the response vertex \(r\)
is not a descendant of \(t\) in this graph. It is also not \(t\).
The pointwise nondescendant assertion of
\cref{thm:p16-causal-intervention-laws}, applied in this modified
model, shows that subsequently fixing \(t=x\) leaves the response
unchanged at every original outcome \(\omega\).

The two fixed interventions are at distinct vertices. Their equation
replacements commute, as proved in the same intervention theorem: in
either order the equation at \(t\) is the constant \(x\), the
equation at \(h\) is the constant \(m\), and every other equation
and outcome weight is unchanged. Hence their unique solutions give
the pointwise identity
\[
 Y^{x,m}(\omega)=Y^m(\omega)
       \qquad\text{for every }\omega\in\Omega.
\]
No observation or positivity hypothesis is used for this identity.

Now start instead with \(\mathcal M_{t=x}\). Its graph is
\(G_{\overline t}\), its solution coordinates include \(M^x,Y^x\),
and its disturbances remain mutually independent. The graph lemma
says that the empty set blocks every back-door path from \(h\)
to \(r\) in this graph. The empty adjustment set also satisfies
the nondescendant condition. Whenever \(\mathbb P(M^x=m)>0\),
the overlap hypothesis of the back-door theorem in this model is
exactly this positive probability, since there is a single empty
conditioning tuple of mass one. Applying
\cref{thm:p16-finite-back-door-adjustment} in \(\mathcal M_{t=x}\)
therefore gives
\[
 \mathbb P(Y^x=y\mid M^x=m)=\mathbb P(Y^{x,m}=y).
\]
The intervention on the right really does replace the mediator's
equation in the already treatment-replaced model; the preceding
commutation argument identifies it with the displayed joint replacement.
The pointwise identity then changes its right side to
\(\mathbb P(Y^m=y)\).

Finally partition the original probability space by the finitely many
values of the random variable \(M^x\). Finite total probability
gives
\[
 \mathbb P(Y^x=y)
 =\sum_{m:\mathbb P(M^x=m)>0}
       \mathbb P(M^x=m)\mathbb P(Y^x=y\mid M^x=m).
\]
Substitute the established identity on each positive stratum. For a
remaining mediator value, \(\mathbb P(Y^m=y)\) is still an ordinary
intervention probability defined by total equation replacement. Its
product with \(\mathbb P(M^x=m)=0\) is thus a defined zero.
Inserting these terms proves the sum over all \(S_h\). This
argument does not assume \(\mathbb P(X=x)>0\) in the original
model, or assign a conditional value on a null event.
\end{proof}
\end{annexitem}

\begin{annexitem}{annex:p16.finite-front-door-identification}{thm:p16-finite-front-door-identification}{The three stages and every observational denominator}
\begin{proof}
\paragraph{Treatment to mediator.}
The second front-door condition says exactly that the empty adjustment
set blocks every back-door path from \(t\) to \(h\). Its
nondescendant requirement is vacuous. The assumed
\(\mathbb P(X=x)>0\) is the overlap condition for this empty
set. The back-door theorem consequently gives, for every \(m\),
\[
 \mathbb P(M^x=m)=\mathbb P(M=m\mid X=x).
\]
The denominator of this conditional ratio is the stated positive
probability of \(X=x\). In particular, the set of mediator values
with positive intervention mass at this \(x\) is exactly
\(\mathcal H_x\).

\paragraph{Mediator to response.}
Fix \(m\in\mathcal H_x\). By the first graph consequence,
\(t\) is a nondescendant of \(h\). It is distinct from the
mediator and response vertices. Together with the third front-door
condition, this says that \(\{t\}\) satisfies the back-door
criterion for intervening at \(h\) and observing \(r\).
For every positive treatment stratum \(x'\in\mathcal X_+\),
the front-door overlap hypothesis supplies
\(\mathbb P(M=m,X=x')>0\). This is precisely the mediator
intervention's overlap condition across the proposed adjustment strata.
The same back-door theorem, now with treatment \(M\) and adjustment
variable \(X\), gives
\[
 \mathbb P(Y^m=y)
   =\sum_{x'\in\mathcal X_+}
      \mathbb P(Y=y\mid M=m,X=x')\mathbb P(X=x').
\]
Every denominator here is the explicitly positive joint probability
\(\mathbb P(M=m,X=x')\). The identity is needed only for
\(m\in\mathcal H_x\); no conditional ratio is assigned for
an unused mediator value with inadequate support.

\paragraph{Composition and identification.}
By \cref{thm:p16-front-door-composition}, the response intervention
law is the sum of \(\mathbb P(M^x=m)\mathbb P(Y^m=y)\).
The first stage identifies the first factor and shows that all terms
outside \(\mathcal H_x\) vanish. The second stage supplies the
second factor for every retained term. Substituting both identities
gives exactly the displayed observational formula.

All probabilities on its right involve only the three coordinates
\(X,M,Y\). They are finite marginal sums and conditional ratios
of the law on any observed set \(O\) containing those coordinates.
If two models in the specified class have equal observed laws, their
sets \(\mathcal X_+\) and \(\mathcal H_x\), all marginal
weights, and all positive-denominator ratios agree. The formula
therefore gives the same intervention mass for every \(y\), which
is identification in \cref{def:p16-causal-identification}.
Equal finite response laws have equal sums against any specified
real-valued response function. Applying the result separately at two
eligible treatment values gives equal intervention contrasts as well.
Neither the latent coordinates' observational law nor a sampling
guarantee has been used.
\end{proof}
\end{annexitem}

\begin{annexitem}{annex:p16.front-door-binary-calculation}{thm:p16-front-door-binary-calculation}{The latent-confounding calculation in full}
\begin{proof}
The outcome weights are nonnegative and sum to
\[
 \left(\tfrac12+\tfrac12\right)
  \left(\tfrac45+\tfrac15\right)
  \left(\tfrac34+\tfrac14\right)=1.
\]
Their product form proves mutual independence of the three coordinate
disturbances. Adding the constant disturbance leaves mutual independence
true, since its sole marginal and joint-compatible value has probability
one. The graph is acyclic with order \(U,X,M,Y\), and all equations
are total on their finite domains.

\paragraph{Every graph condition.}
The only directed path from \(X\) to \(Y\) is
\(X\to M\to Y\), which contains the mediator. The only back-door
path from \(X\) to \(M\) is \(X\leftarrow U\to Y\leftarrow M\).
It has a collider at \(Y\), and is blocked by the empty set.
The only back-door path from \(M\) to \(Y\) is
\(M\leftarrow X\leftarrow U\to Y\); its internal vertex \(X\)
is a noncollider, so conditioning on \(X\) blocks it. These are
all such paths: the underlying four-vertex graph is the cycle with
undirected edges \(\{U,X\}\), \(\{X,M\}\), \(\{M,Y\}\),
\(\{Y,U\}\), so each of these
pairs has exactly the two simple paths just described, one causal
and one back-door. This proves all three graph conditions.

\paragraph{The eight full assignments.}
For given \(u,x,m\), the unique disturbance coordinates are
\(a=x\oplus u\) and \(b=m\oplus x\). The response is one
exactly when \(u=1\) or \(m=1\). Thus direct evaluation gives
the following full law; every unlisted assignment has mass zero:
\[
\begin{array}{cccc|c}
 U&X&M&Y&\text{mass}\\\hline
 0&0&0&0&12/40\\
 0&0&1&1&4/40\\
 0&1&0&0&1/40\\
 0&1&1&1&3/40\\
 1&0&0&1&3/40\\
 1&0&1&1&1/40\\
 1&1&0&1&4/40\\
 1&1&1&1&12/40
\end{array}
\]
For example, the four possible products
\(\frac12\alpha(a)\beta(b)\), in the order
\((a,b)=(0,0),(0,1),(1,0),(1,1)\), are
\(12/40,4/40,3/40,1/40\); the table uses exactly the indicated
inverse coordinates at each row. Its masses sum to \(40/40\).
Grouping rows with the same \((X,M,Y)\) gives the body table.

Both treatment values have mass \(20/40=1/2\). The positive
\((X,M)\) masses in the order \((0,0),(0,1),(1,0),(1,1)\)
are \(15/40,5/40,5/40,15/40\). Therefore both mediator values
occur at every treatment value, proving all the required overlap.
The conditional mediator law is
\[
\begin{array}{c|cc}
 &M=0&M=1\\\hline
 X=0&3/4&1/4\\
 X=1&1/4&3/4
\end{array}
\]
and the conditional response probabilities are
\[
\begin{array}{c|cc}
 &X=0&X=1\\\hline
 M=0&1/5&4/5\\
 M=1&1&1
\end{array}
\]
Indeed the first row is \((3/40)/(15/40)=1/5\) and
\((4/40)/(5/40)=4/5\), while at \(M=1\) every response is one.
All four denominators are positive.

\paragraph{The observational and interventional responses.}
Summing the body table over the mediator gives response-one mass
\(8/40\) in the \(X=0\) stratum and \(19/40\) in the
\(X=1\) stratum. Dividing by \(1/2\) gives \(2/5\) and
\(19/20\); their difference is \(11/20\).
For the inner front-door sum, averaging each row of the response
table over \(\mathbb P(X=0)=\mathbb P(X=1)=1/2\) gives
\[
 q_0=\tfrac12(\tfrac15+\tfrac45)=\tfrac12,
 \qquad q_1=\tfrac12(1+1)=1.
\]
Here \(q_m\) is the computed response-one intervention probability
at mediator value \(m\), not the outcome-space mass \(q(u,a,b)\).
The outer averages are
\[
 \tfrac34 q_0+\tfrac14 q_1=\tfrac58,\qquad
 \tfrac14 q_0+\tfrac34 q_1=\tfrac78.
\]

Directly fixing \(X=x\) gives \(M=x\oplus b\), while the
root \(U=u\) remains fair and independent of \(b\). The equation
for \(Y\) is unchanged. Its response-one event is the disjoint
union of \(\{U=1\}\) and \(\{U=0,M^x=1\}\), so
\[
 \mathbb P(Y^x=1)
   =\tfrac12+\tfrac12\mathbb P(x\oplus U_M=1).
\]
At \(x=0\) the final probability is \(1/4\), and at \(x=1\)
it is \(3/4\), giving the same \(5/8,7/8\) values. Their
difference is \(1/4\). As a bit-valued response has expectation
equal to its probability of being one, these are also the stated
mean contrasts.

Finally, the only other measured vertex eligible to be considered
for adjustment of \(X\) and \(Y\) is \(M\). The empty set
does not block \(X\leftarrow U\to Y\). The set \(\{M\}\)
violates the nondescendant condition, because \(X\to M\).
Thus no subset of the other measured coordinates satisfies that
back-door criterion, although the front-door formula identifies the
effect. This is a comparison of the two proved sufficient criteria,
not a claim that failure of back-door adjustment implies
nonidentifiability.
\end{proof}
\end{annexitem}

\begin{annexitem}{annex:p16.front-door-bypass-counterexample}{thm:p16-front-door-bypass-counterexample}{The formula fails while the effect remains identified}
\begin{proof}
The four exogenous weights \(\frac12\beta(b)\), indexed by a
fair bit for \(X\) and the disturbance bit \(b\) for \(M\),
sum to one and are products. The two coordinate disturbances and
the constant response disturbance are mutually independent. The
topological order \(X,M,Y\) evaluates the total equations to
\(X=c,M=c\oplus b,Y=\mathbf1_{\{c=1\text{ or }c\oplus b=1\}}\).
Its complete observational law is
\[
\begin{array}{ccc|c}
 X&M&Y&\text{mass}\\\hline
 0&0&0&3/8\\
 0&1&1&1/8\\
 1&0&1&1/8\\
 1&1&1&3/8
\end{array}
\]
All other assignments have mass zero, and the displayed masses sum
to one. Both treatment values have probability \(1/2\), and all
four treatment-mediator pairs have positive probability.

The vertex \(X\) has no incoming edge, so it has no back-door
path to \(M\), proving the second front-door condition. The only
back-door path from \(M\) to \(Y\) is \(M\leftarrow X\to Y\),
blocked by its noncollider \(X\), proving the third condition.
The direct edge \(X\to Y\) is a directed path that avoids
\(M\), so the first condition fails.

The conditional mediator probabilities are \((3/4,1/4)\) at
\(X=0\) and \((1/4,3/4)\) at \(X=1\). From the deterministic
response equation,
\[
 \mathbb P(Y=1\mid M=0,X=0)=0,\quad
 \mathbb P(Y=1\mid M=0,X=1)=1,
\]
and \(\mathbb P(Y=1\mid M=1,X=x')=1\) for both \(x'\).
Thus the inner sum of the front-door expression is \(1/2\) at
\(M=0\) and one at \(M=1\). Its outer sums are \(5/8\) and
\(7/8\), just as stated.

Under equation replacement \(X=0\), the mediator equals its
disturbance bit and the response equals that mediator, so
\(\mathbb P(Y^{X=0}=1)=1/4\). Under \(X=1\), the response
equation is identically one, so \(\mathbb P(Y^{X=1}=1)=1\).
The proposed front-door functional is therefore wrong at both values.

There is nevertheless no back-door path from the root \(X\) to
\(Y\), and \(\mathbb P(X=x)>0\) for both \(x\). The
empty-set back-door theorem applies to this treatment and response,
identifying its effect as \(\mathbb P(Y=y\mid X=x)\). The
observational table gives response-one conditionals \(1/4\) and
one, agreeing with direct intervention. Thus an observed-law functional
does identify the effect in this graph class. The failed front-door
expression supplies no counterexample to that identification claim.
\end{proof}
\end{annexitem}
